\section*{الفصل \ref{ch:b2:atmospheric-thermodynamics} --- الديناميكا الحرارية للغلاف الجوي}

\begin{solution}{exo:b2:atmospheric-thermodynamics:1}
(أ) $e = e_s(\qty{278.15}{K}) = \qty{8.73}{hPa}$ و$e_s(\qty{283.15}{K}) = \qty{12.32}{hPa}$.
(ب) $8.73/12.32 = \qty{70.9}{\percent}$. (ج) $r = 0.622 \times 8.73/(850 - 8.73) =
\qty{6.45}{g/kg}$ و$q = r/(1 + r) = \qty{6.41}{g/kg}$.
\end{solution}

\begin{solution}{exo:b2:atmospheric-thermodynamics:2}
(أ) $\mathrm{d}\ln e_s/\mathrm{d}T = 5419/T^2$: \qty{7.26}{\percent} لكل كلفن عند
\qty{273.15}{K}، و\qty{6.31}{\percent} عند \qty{293.15}{K}، و\qty{5.90}{\percent} عند
\qty{303.15}{K}. (ب) $e_s(\qty{20}{\celsius}) = \qty{23.66}{hPa}$ و$e_s(\qty{25}{\celsius})
= \qty{32.26}{hPa}$؛ والنسبة 1.363، مقابل $1.063^5 = 1.357$. فالمعدّل مشتقّ
$\ln e_s$، فيكون العامل على مدى \qty{5}{K} هو $\exp(5 \times 0.0620) = 1.363$، حيث 0.0620
المعدّل عند نقطة المنتصف \qty{22.5}{\celsius}؛ أمّا تركيب \qty{6.3}{\percent} خمس مرّات فقريب
لكنه غير دقيق. (ج) $23.66/23.39 = 1.012$: فالصيغة أعلى ممّا ينبغي بمقدار \qty{1.2}{\percent}، لأنها
تحتفظ بالحرارة الكامنة للنقطة الثلاثية، في حين تتناقص الحرارة الكامنة كلّما ارتفعت درجة
الحرارة؛ فميل الصيغة أشدّ ممّا ينبغي، ويكبر خطؤها مع البعد عن
\qty{0}{\celsius}.
\end{solution}

\begin{solution}{exo:b2:atmospheric-thermodynamics:3}
(أ) $125 \times (30 - 14) = \qty{2000}{m}$. (ب) $30 - 9.8 \times 2.0 = \qty{10.4}{\celsius}$.
(ج) $14 - 1.8 \times 2.0 = \qty{10.4}{\celsius}$: تلتقي درجة الحرارة ونقطة الندى، كما يجب عند
مستوى التشبّع؛ والقاعدة هي هذا الحساب. (د) $125 \times (33 - 14) = \qty{2375}{m}$:
ترتفع القواعد كلّما سخّن العصرُ هواءَ الأرض.
\end{solution}

\begin{solution}{exo:b2:atmospheric-thermodynamics:4}
(أ) يقع منحنى نقطة الندى على منحنى الحرارة من الأرض حتى نحو \qty{908}{hPa}:
فنصف الكيلومتر الأدنى مشبع، أو قريب من التشبّع. (ب) بين نحو 908 و\qty{889}{hPa}
ينعطف منحنى الحرارة نحو اليمين، إلى قيم أدفأ، مع تناقص الضغط: أي
ترتفع الحرارة مع الارتفاع. (ج) \omterm{def:b2:atmospheric-thermodynamics:lcl}{مستوى التكاثف بالرفع} عند نحو \qty{964}{hPa}، على ارتفاع بضع عشرات الأمتار؛ \omterm{def:b2:atmospheric-thermodynamics:parcel-levels}{ومستوى الحمل الحرّ} عند
نحو \qty{730}{hPa}؛ \omterm{def:b2:atmospheric-thermodynamics:parcel-levels}{ومستوى التوازن} عند نحو \qty{235}{hPa}. (د) عند قمّة الانعكاس تكون نقطة الندى
أدنى من الحرارة بمقدار \qty{11}{K}، وأدنى منها بمقدار \qty{20}{K} عند \qty{850}{hPa}: فهواء دافئ شديد الجفاف
يعلو هواءً باردًا رطبًا. ورطوبته النسبية منخفضة، فلم يكتسب
دفأه بالتكاثف؛ والهواء الهابط يسخن بالانضغاط بالمعدّل الأديباتي الجاف مع
احتفاظه بنسبة خلطه، وهذا ما يعطي مثل هذه الطبقة.
\end{solution}

\begin{solution}{exo:b2:atmospheric-thermodynamics:5}
(أ) $T_v = 303.15 \times (1 + 0.020/0.622)/1.020 = \qty{306.76}{K}$؛ ويعطي التقريب
$303.15 \times (1 + 0.61 \times 0.020) = \qty{306.85}{K}$. (ب) $\rho = p/(R_dT_v) = 10^5/(287.0
\times 306.76) = \qty{1.1358}{kg/m^3}$، مقابل $10^5/(287.0 \times 303.15) =
\qty{1.1494}{kg/m^3}$ للهواء الجاف: أي أخفّ بنسبة \qty{1.2}{\percent}. (ج) بمقدار $T_v - T =
\qty{3.6}{K}$. (د) لا يحتاج الطرد الصاعد إلا إلى أن يكون أخفّ من محيطه: وفوق البحار الدافئة
تضيف رطوبة الهواء طفوًا يعادل عدّة كلفنات، أي ما يعادل زيادة
معتبرة في الحرارة.
\end{solution}

\begin{solution}{exo:b2:atmospheric-thermodynamics:6}
(أ) $e_s$ = 6.11 و23.66 و\qty{43.54}{hPa}؛ و$r_s = 0.622\,e_s/(1000 - e_s)$ = 3.82 و15.07
و\qty{28.31}{g/kg}؛ و$\Gamma_s$ = 6.48 و4.20 و\qty{3.44}{K/km}. (ب) يبرد الطرد الجاف بمقدار
\qty{9.8}{K}، والمشبع بمقدار \qty{4.2}{K}: فالطرد المشبع أدفأ بمقدار \qty{5.6}{K}.
(ج) في الأعالي يكون الهواء أبرد، و$r_s$ أصغر، ولا يبقى إلا قليل من البخار ليتكاثف،
ويعود $\Gamma_s$ نحو $\Gamma_d$.
\end{solution}

\begin{solution}{exo:b2:atmospheric-thermodynamics:7}
(أ) $e = e_s(\qty{291.65}{K}) = \qty{21.52}{hPa}$، و$r = 0.622 \times 21.52/(971.1 - 21.52) =
\qty{14.1}{g/kg}$؛ و$\theta = 292.15 \times (1000/971.1)^{0.2857} = \qty{294.6}{K}$؛ و$\theta_e =
294.6 \exp[2.501 \times 10^6 \times 0.0141/(1004.5 \times 291.5)] = 294.6 \times 1.128 =
\qty{332.3}{K}$. (ب) $e_s(\qty{264.25}{K}) = \qty{3.13}{hPa}$، و$r_s = 0.622 \times 3.13/(500 -
3.13) = \qty{3.92}{g/kg}$؛ و$\theta = 264.25 \times 2^{0.2857} = \qty{322.1}{K}$؛ والطرد
مشبع، فيكون $T_{\text{LCL}} = T$ و$\theta_e = 322.1 \exp[2.501 \times 10^6 \times
0.00392/(1004.5 \times 264.25)] = \qty{334.3}{K}$. (ج) ارتفعت $\theta$ بمقدار \qty{27.5}{K}، وقد سخّنتها
الحرارة الكامنة لمقدار \qty{10.2}{g/kg} من البخار الذي تكاثف؛ أمّا $\theta_e$ فتغيّرت بنسبة
\qty{0.6}{\percent}، ضمن دقّة صيغتها التقريبية: فهي محفوظة.
\end{solution}

\begin{solution}{exo:b2:atmospheric-thermodynamics:8}
(أ) $\text{CAPE} = 287.0 \times 4 \times \ln(700/300) = 287.0 \times 4 \times 0.847 =
\qty{973}{J/kg}$. (ب) $\sqrt{2 \times 973} = \qty{44.1}{m/s}$. (ج) بين 1
و\qty{1500}{J/kg}: موجبة، لكنها ليست كبيرة. (د) شريط ذو عرض أفقي ثابت بين
منحنيين؛ ووفق \cref{prop:b2:atmospheric-thermodynamics:areas} تتناسب مساحته مع $\int
\Delta T\,\mathrm{d}\ln p = 4 \ln(700/300)$، أيًّا كان ميل خطوط تساوي الحرارة.
\end{solution}

\begin{solution}{exo:b2:atmospheric-thermodynamics:9}
(أ) $D_c = \sqrt{3 \times 0.61 \times 80^3/2.1}\,\unit{nm} = \qty{668}{nm} =
\qty{0.668}{\micro m}$؛ و$s_c = 2A/(3D_c) = 2 \times 2.1/(3 \times 668) = 0.00210 =
\qty{0.210}{\percent}$. (ب) بحلّ $s_c = \sqrt{4A^3/(27\kappa D_d^3)}$ بالنسبة إلى $D_d$: $D_d =
\bigl(4A^3/(27\kappa s_c^2)\bigr)^{1/3}$، أي \qty{63}{nm} عند \qty{0.3}{\percent} و\qty{131}{nm} عند
\qty{0.1}{\percent}. (ج) كلّ جسيم بين 63 و\qty{131}{nm} يتنشّط الآن أيضًا: قطيرات
أكثر، وأصغر مع الكمّية نفسها من الماء السائل.
\end{solution}

\begin{solution}{exo:b2:atmospheric-thermodynamics:10}
(أ) $w = \sqrt{2 \times 130} = \qty{16.1}{m/s}$، وهذا مستحيل لهواء ساكن قرب الأرض.
(ب) درجة الحرارة الكامنة لقمّة الانعكاس هي $293.65 \times (1000/889.1)^{0.2857} =
\qty{303.7}{K}$؛ وطرد سطحي بهذه $\theta$ له $T = 303.7 \times (971.1/1000)^{0.2857}
= \qty{301.2}{K}$، أي \qty{28.0}{\celsius}. (ج) $125 \times (28.0 - 18.5) = \qty{1190}{m}$ فوق
الأرض، فوق قمّة الانعكاس القديمة. (د) جبهة، ترفع الهواء على منحدرها؛ والصعود البطيء
للهواء أمام اضطراب في الأعالي، الذي يبرّد الغطاء نفسه.
\end{solution}

\begin{solution}{exo:b2:atmospheric-thermodynamics:11}
(أ) حين $r_s \to 0$، يؤول البسط إلى $g$ والمقام إلى $c_p$. (ب) اكتب
$L^2r_s/(R_vT^2) = c_p \times \bigl(Lr_s/(R_dT)\bigr) \times LR_d/(R_vTc_p)$، مع
$LR_d/(R_vTc_p) = 5.1$ عند \qty{303}{K}: فحين $Lr_s/(R_dT) \gg 1$، يكون أيضًا $L^2r_s/(R_vT^2) \gg c_p$،
و$\Gamma_s \to g\,(Lr_s/(R_dT))/(L^2r_s/(R_vT^2)) = gR_vT/(LR_d) = 9.81 \times 461.5 \times
303/(2.501 \times 10^6 \times 287.0) = \qty{1.91}{K/km}$. (ج) مع $T_d = T$، يكون
$\Gamma_{\text{ندى}} = gR_vT/(LR_d)$، وهو نفسه. ففي تلك النهاية تكون الحرارة الكامنة كبيرة إلى حدّ أن
الطرد لا يستطيع تكثيف شيء يُذكر من بخاره دون أن يسخن كثيرًا فوق نقطة
تشبّعه: فيحتفظ بنسبة خلطه ويتبع خطّ \omterm{def:b2:atmospheric-thermodynamics:moist-air}{نسبة الخلط} الثابتة، الذي تنخفض على امتداده
الحرارة بالمعدّل $\Gamma_{\text{ندى}}$. (د) عند \qty{30}{\celsius} و\qty{1000}{hPa}،
$Lr_s/(R_dT) = 0.81$، ليس أكبر بكثير من 1، و$L^2r_s/(R_vT^2) = 4180$، أي أربعة أضعاف $c_p$:
$\Gamma_s = \qty{3.4}{K/km}$، لا يزال بعيدًا عن \qty{1.9}{K/km}. فالنهاية لا تُبلغ.
\end{solution}

\begin{solution}{exo:b2:atmospheric-thermodynamics:12}
(أ) $2.868/2.599 = 1.104$ و$1.257/1.032 = 1.218$: أي \qty{10.4}{\percent}
و\qty{21.8}{\percent}. (ب) تتنشّط القطيرات عند بضعة أعشار الواحد في المئة من \omterm{def:b2:atmospheric-thermodynamics:microphysics}{فوق التشبّع}،
أي أقلّ بعشرات المرّات. (ج) تنمو البلّورات سريعًا بالترسّب البخاري؛ فتأخذ البخار من الهواء،
فينخفض تحت التشبّع فوق القطيرات؛ فتتبخّر القطيرات ويذهب بخارها إلى
الجليد (عملية برجيرون)، حتى تتجلّد السحابة. (د) عند \qty{0}{\celsius} يتساوى ضغطا
البخار المشبع: فهما يلتقيان عند النقطة الثلاثية، وتؤول النسبة إلى الواحد
قربها.
\end{solution}

\begin{solution}{pb:b2:atmospheric-thermodynamics:1}
\textbf{1.} $e = e_s(\qty{291.65}{K}) = \qty{21.5}{hPa}$.
\textbf{2.} $r = 0.622 \times 21.5/(971.1 - 21.5) = \qty{14.1}{g/kg}$.
\textbf{3.} $q = 14.1/1.0141 = \qty{13.9}{g/kg}$.
\textbf{4.} $e_s(\qty{292.15}{K}) = \qty{22.2}{hPa}$؛ و$21.5/22.2 = \qty{97}{\percent}$.
\textbf{5.} $0.61 \times 0.0141 \times 292.15 = \qty{2.5}{K}$.
\textbf{6.} $125 \times 0.5 = \qty{63}{m}$؛ و$p = 971.1 \exp[-9.81 \times 62.5/(287.0 \times 294.7)]
= \qty{964}{hPa}$.
\textbf{7.} $\theta = 292.15 \times (1000/971.1)^{0.2857} = \qty{294.6}{K}$.
\textbf{8.} $16.9 - 17.3 = \qty{-0.4}{K}$: أبرد.
\textbf{9.} $(20.5 - 16.4)/(1102 - 924) = 4.1/178$: ترتفع الحرارة بمعدّل \qty{23}{K/km}، أي
معدّل انخفاض قدره \qty{-23}{K/km}. إنه انعكاس حراري: مستقرّ استقرارًا مطلقًا، غطاء.
\textbf{10.} $15.4 - 20.5 = \qty{-5.1}{K}$.
\textbf{11.} $e_s(\qty{-1.8}{\celsius}) = \qty{5.4}{hPa}$ و$e_s(\qty{18.6}{\celsius}) =
\qty{21.7}{hPa}$: أي \qty{25}{\percent}. وهو دافئ وجاف، فقد يكون هواءً هبط وسخن
بالانضغاط، أو هواءً سُخّن فوق أرض جافّة في اتجاه منبع الريح؛ وليس هو الهواء الرطب الذي تحته.
\textbf{12.} $18.6 + 18.4 = \qty{37.0}{K}$ على مدى $5725 - 1489 = \qty{4236}{m}$، أي \qty{8.73}{K/km}: بين $\Gamma_s$
و$\Gamma_d$، فالطبقة غير مستقرّة استقرارًا مشروطًا.
\textbf{13.} $-18.4 - (-8.9) = \qty{-9.5}{K}$: عدم استقرار شديد (من $-8$ إلى $-10$).
\textbf{14.} من الأرض صعودًا: 0، $-0.4$، $-0.2$، $-5.1$، $-4.9$، $+1.6$، $+9.5$، $+10.0$،
$+5.6$، $+1.5$، $-9.5$~\unit{K}.
\textbf{15.} بين 850 ($-4.9$) و\qty{700}{hPa} ($+1.6$)، عند الكسر $4.9/6.5 = 0.754$:
$\ln p = \ln 850 + 0.754 \ln(700/850)$، و$p = \qty{734}{hPa}$؛ و$z = 1489 + 0.754 \times 1621 =
\qty{2711}{m}$، أي \qty{2.37}{km} فوق الأرض.
\textbf{16.} بين 250 ($+1.5$) و\qty{200}{hPa} ($-9.5$)، عند الكسر $1.5/11.0 =
0.136$: $p = \qty{242.5}{hPa}$ و$z = 10560 + 0.136 \times 1446 = \qty{10757}{m}$، أي نحو
\qty{200}{m} تحت التروبوبوز المسجَّل (عند \qty{10954}{m}).
\textbf{17.} المقاطع (متوسط $\Delta T$ مضروبًا في $\ln$ نسبة الضغطين): $0.80 \times 0.0474
= 0.038$؛ و$5.55 \times 0.3365 = 1.867$؛ و$9.75 \times 0.2231 = 2.176$؛ و$7.80 \times 0.2877 =
2.244$؛ و$3.55 \times 0.1823 = 0.647$؛ و$0.75 \times 0.0305 = 0.023$؛ والمجموع \qty{6.995}{K}.
$\text{CAPE} = 287.0 \times 6.995 = \qty{2010}{J/kg}$، إلى ثلاثة أرقام.
\textbf{18.} بين 1500 و\qty{2500}{J/kg}: كبيرة.
\textbf{19.} $-0.2 \times 0.0486 - 0.3 \times 0.0189 - 2.65 \times 0.0207 - 5.0 \times 0.0450 -
2.45 \times 0.1467 = -0.655$~\unit{K}؛ و$\text{CIN} = 287.0 \times (-0.655) = \qty{-188}{J/kg}$.
\textbf{20.} لا يحتفظ الجدول إلا بأحد عشر مستوى: وبينها تستبدل قاعدة شبه المنحرف بالمنحنيات
قطعًا مستقيمة بدلالة $\ln p$ فتفوتها انحناءاتها. أمّا الشبكة ذات الخطوة \qty{1}{hPa} فتتبع المنحنيات. وتضيف الحرارة الافتراضية
طفو بخار الماء في الطرد، وهو أوفر في الطرد المشبع منه في
الهواء الجاف المحيط به: فتكبر CAPE ويصغر \omterm{def:b2:atmospheric-thermodynamics:cape}{كبح الحمل}.
\textbf{21.} $\sqrt{2 \times 2010} = \qty{63}{m/s}$.
\textbf{22.} $10757 - 2711 = \qty{8046}{m}$ بسرعة $63.4/2 = \qty{31.7}{m/s}$: أي \qty{254}{s}، نحو أربع
دقائق.
\textbf{23.} وزن الماء المتكاثف الذي يحمله التيّار الصاعد؛ واختلاطه بالهواء الأجفّ
المحيط به (الإدخال)، الذي يبرّده إذ تتبخّر القطيرات.
\textbf{24.} تعطي نظرية الطرد التيّارَ الصاعد لعاصفة يغذّيها هذا الهواء سرعةً تقارب
\textbf{\qty[text-series-to-math=true]{63}{m.s^{-1}}}، وهي حدّ أعلى.
\end{solution}
